\documentclass{handout}

\assignment{Homework 8}
\title{Circuits}
\duedate{Wednesday November 12, 2025}

\begin{document}
\maketitle

\hwinstructions

\begin{questions}
    \question \textbf{Simple circuit:} 	Consider the circuit shown below. The terminal voltage of the battery is  18.00 V.
    \begin{parts}
        \part Find the equivalent resistance of the circuit.
        \part Find the current through each resistor.
        \part Find the potential drop across each resistor.
        \part Find the power dissipated by each resistor.
        \part Find the power supplied by the battery.
    \end{parts}
    \begin{figure}[h!]
        \centering
        \includegraphics[height=1.5in]{Images/Simple circuit.pdf}
    \end{figure}

    \question \textbf{Christmas light wattage:} Christmas lights are usually wired as follows: groups of 50 are wired in \textbf{series}, then several of these groups are wired in \textbf{parallel}. A common number to have is a 300 light strand, so there are 6 such groups.
    \begin{parts}
        \part If $R$ is the resistance of one light bulb while operating, write an expression for the resistance of one group of 50 bulbs in series.
        \part Write an expression for the resistance of 6 groups wired in parallel
        \part A strand of 300 LED Christmas lights draws 21 W, whereas a strand of incandescent Christmas lights draws 72 W. What are the total resistances of each whole string of lights if plugged into a 120 V outlet?
        \part What is the resistance of each individual bulb, incandescent and LED?
        \part Circuit breakers trip when the \textbf{current} to a group of outlets exceeds some set amount, usually 15 A. If you wanted to stay within 80\% of this current value, how many LED strands could you connect end-to-end (these are being wired in \textbf{parallel})? How many incandescent strands?
    \end{parts}

    \begin{figure}[h!]
        \centering
        \includegraphics[width=0.75\linewidth]{Images/christmas lights.pdf}
    \end{figure}

    \question \textbf{Kettle power}: It takes 334 kJ of heat energy to boil water. You would like to design an electric kettle that boils water in 3 minutes.
    \begin{parts}
        \part How much power, in Watts (W), should your electric kettle put out to accomplish this task?
        \part How much current would this require in an American home, where the voltage is 120 V?
        \part What should be the resistance of the heating element that you put in the kettle?
        \part Nichrome wire has a conductivity of $\sigma =10^6 (\Omega \ \text{m})^{-1}$. You have some 24 gauge nichrome wire, which has a cross section of $0.25$ mm$^2$. How long of a nichrome wire will you need to use to get the resistaance you found in part c? (This element is usually coiled so you can fit it in a smaller space).
    \end{parts}

    \question \textbf{Kirchoff's rules:} Given the circuit below, find $V_1,I_2,I_3$.
    \begin{parts}
        \part Using Kirchoff's junction rule, write down an equation for the junction shown.
        \part Write down two loop rule equations.
        \part Solve this system of equations for the three unknowns. Show your work.
    \end{parts}
    \begin{figure}[h!]
        \centering
        \includegraphics[width=0.5\linewidth]{Images/kirchoffCircuit.png}
    \end{figure}
\end{questions}

\end{document}
